% =============================================================================
% sections/design/drone/budgets/power.tex  -  Drone Power Budget
% Teams: drone (best guess - unconfirmed) / Status: EXAMPLE (scope approximation)
%
% EXAMPLE CONTENT for scope/layout approximation - not final report text.
% Adapted from the Preliminary Report's "Drone Power Budget" subsection.
% Marked with \exampleflag on its heading in the assembler.
% =============================================================================

The power budget is evaluated for a single operating condition: \emph{hover}
(propulsion with four rotors, attitude control, and baseline electronics),
$P_{\text{hov}} = P_{\text{prop}} + P_{\text{control,avg}}$. Propulsion
power is estimated from induced and profile power (momentum theory adapted
for multicopter hover). For a quadcopter ($N_R = 4$), the total weight is
shared equally among the rotors:
\[
    T_{\text{total}} = mg, \qquad
    T_{\text{prop}} = \frac{T_{\text{total}}}{4}, \qquad
    \sigma = \frac{Nc}{\pi R}, \qquad
    P_{i,h} = k\,\frac{T_{\text{prop}}^{3/2}}{\sqrt{2\rho A}},
\]
\[
    P_0 = \rho A\, V_{\text{tip}}^3\, \frac{\sigma c_{d0}}{8}, \qquad
    P_{\text{mr,per}} = P_{i,h} + P_0, \qquad
    P_{\text{mr}} = 4\,P_{\text{mr,per}}, \qquad
    P_{\text{prop}} = P_{\text{mr}}.
\]

Assumptions: $m = \qty{2.75}{\kilo\gram}$ (total vehicle mass,
\autoref{tab:drone_mass_budget}), $\rho = \qty{1.225}{\kilo\gram\per\meter\cubed}$
(sea-level air density), $R = \qty{0.127}{\meter}$ (rotor radius per
propeller), $A = \pi R^2 = \qty{0.0506}{\meter\squared}$ (disk area per
rotor), $V_{\text{tip}} = \qty{118.1}{\meter\per\second}$ (tip speed,
assumed \qty{11100}{\rpm}), $c_{d0} = 0.025$ (blade profile drag
coefficient), $k = 1.1$ (induced-power correction factor, from a range of
1.10--1.15).

\begin{table}[!htbp]
    \centering
    \caption{Drone hover power budget}
    \label{tab:drone_power_budget}
    \begin{tabular}{L{5cm} L{4.5cm} R{2.5cm}}
        \toprule
        \textbf{Term} & \textbf{Expression / basis} & \textbf{Value} \\
        \midrule
        Induced power (per motor)   & $P_{i,h}$                                  & \qty{54.7}{\watt} \\
        Rotor solidity $\sigma$     & $Nc/(\pi R)$                                & $0.0595$ \\
        Profile power (per motor)   & $P_0$                                       & \qty{18.99}{\watt} \\
        Mechanical rotor power      & $P_{\text{mr}} = 4(P_{i,h}+P_0)$           & \qty{294.76}{\watt} \\
        Propulsion power            & $P_{\text{prop}} = P_{\text{mr}}$          & \qty{294.76}{\watt} \\
        Control average             & $P_{\text{control,peak}} D$, $D \in [0.05,0.15]$ & \qty{20}{\watt} (nominal) \\
        \textbf{Total hover power}  & $P_{\text{prop}} + P_{\text{control,avg}}$ & \textbf{\qty{314.76}{\watt}} \\
        Hover energy (\qty{60}{\minute}) & $P_{\text{hov}} \cdot \qty{1}{\hour}$  & \qty{314.76}{\watt\hour} \\
        \bottomrule
    \end{tabular}
\end{table}

The Drone is powered by a 5S LiPo battery, sized to deliver the voltage and
energy capacity required to complete the mission.

\todo[inline]{Confirm the battery capacity actually selected covers the
\qty{314.76}{\watt\hour} hourly hover energy with margin, and state that
capacity explicitly (the Preliminary specified "a 5S LiPo battery" without
a capacity figure).}
